% kbconfig.tex                               An idea of Bernard Gaulle 07/29/92
%                                                                         V1.16
%                           Copyright GUTenberg-Bernard Gaulle as in french.doc
%
% Is supposed to run with \catcode`@=11.
%
% mods: Let \csubinverse be known of the rest of the world.       --bg 04/20/93
%       Let 7-bit accent macro produce 8-bit chars when possible. --bg 04/28/93
%       Accent macros must produce only one token  otherwise
%           upper/lower-casing isn't correctly
%           converted (bug found by Jean-Claude Vivalda).         --bg 05/06/93
% V1.3  Define << and >> for any 8-bit usage(pointed by H.Soulard)--bg 08/10/93
% V1.4  Define a \global \FontEncoding as "CMR" or "ECM" to pass
%           to any run-time style such as french.sty  (but let
%           it asis if already defined).                          --bg 08/18/93
%       Add code for kb2lex interface (only for TeX V3.xx)        --bg 08/27/93
%       #4 param of \Madef wasn't a balanced text.                --bg 08/30/93
%       Replace \FontEncoding by \encodingdefault of NFSS2.       --bg 08/31/93
%       Current encoding used if any (\f@encoding).               --bg 09/01/93
%       \@tempa no more redefined bec. used in \input def of NFSS2--bg 09/02/93
%       But ngow using \@tempc...
%       Introduce \protect outside kb2lex because a macro such
%       as \copyright is not expandable due to alignement macros. --bg 09/03/93
%       LaTeX loop over \@tempc solved by \edef to null in[ macros
%       Issue warning when TeX V3 and CMR only.                   --bg 11/22/93
%       \f@encoding may be undefined.                             --bg 12/03/93
%       Allow the NFSS2 user to change of font encoding on the fly--bg 12/15/93
%       \input replaced by \@@input if defined.                   --bg 12/17/93
% V1.5  Accent macros redefined at end to usual if outside MlTeX.
%       \aa and \AA were not redefined for 8bits.                 --bg 12/21/93
% V1.6  Let \ifECM be globaly false if undefined. (pb of Goussard)--bg 03/23/94
% V1.7  Correction of previous change that was wrongly coded.     --bg 03/31/94
% V1.8  \ISOfalse disactivated too much active chars.             --bg 05/09/94
% V1.9  Reduced usage of backslash in lex code.(pointed by Vial). --bg 05/30/94
% V1.10 Give a global default definition for < and >.             --bg 06/07/94
% V1.11 Reset invalid LaTeX2e chars to internal \catcode 11.      --bg 09/27/94
%       Hack the format dumping time to set catcode correctly
%       (normally to remove for the 94-95 winter release of LaTeX)--bg 09/29/94
%       Take care of \I & \SS already defined in 2e.
%       Correct 2e bug regarding OT1 uppercasing of ^^Y, ^^Z      --bg 09/30/94
%       Document these hacks via \wlog.                           --bg 10/03/94
%       Create a latex209.cfg to patch latex209.def which wrongly
%         reset \catcodes.                                        --bg 10/06/94
%       Let 8bit << and >> be equal to \fguillemets and \endfguillemets
%         avoiding them to become back 7bit when a solution will
%         be found...                                             --bg 10/10/94
%       Don't define globally current 2e TextSymbols.
%       Don't redefine accent macros for 2e T1 which includes something
%         like our MlTeX \csubinverse mechanism.                  --bg 10/26/94
%       Rename \: in \r wich is what 2e is using.                 --bg 11/16/94
%       Stronger error messages when running initex TeX+cmr.
%       Define MlTeX \OT1\accent macros for 2e.                   --bg 11/30/94
%       Take in account case when hexa code are undefined for
%         double chars outside MlTeX.                             --bg 12/17/94
%       Set low invalid chars to 11 even if gobbled.
%       Reload of kbconfig needed only for OT1.                   --bg 12/18/94
%       8bit chars in T1 were not defined as accent on letter.    --bg 12/20/94
% V1.13 oe, ae, etc were not activated in T1 ISO.
%       Uppercase of oe, ae, etc were not set to 11 in T1.
%       Force T1 if necessary when emulating 209.                 --bg 01/05/95
% V1.14 \@c@@ wasn't correctly defined for OT1 in 2e.             --bg 01/09/95
%       Avoid hexa codes to point back elsewhere in non-ISO.
%       Enforce \ifHB and encodingdefault wishes.                 --bg 01/10/95
%       Use ^^9c as a gateway to uppercase estzet.
%       \csubinverse didn't use hexa code for non-ISO.
%       \r was undefined in plain TeX.
%       \lccode of \OE was undefined in (and for) plain TeX.
%       Plain def of pounds must not use \char bu \string         --bg 01/14/95
% V1.15 Let \i and \j be i and j inside 2e.                       --bg 02/05/95
% V1.16 Low chars need to be 11, even in 2e.                      --bg 03/06/95
%       Force T1 in latex209.cfg for non-MlTeX engines to ensure
%         that hyphenation be clean. Sugg. by D.B.Roegel.
%       No lcc/ucc for normal T1 ISO slots (correctly set in font).
%       In T1 ISO, chars are not initiated. This was true except
%         for MACRO & LIGS (chars were already activated).
%       \DOhex provides chardef adjustement in hex. True but not
%         for LEX.
%       Plain has an incompatible \i definition for \csubinverse
%        and no definition of \pounds (was wrongly removed by me
%        as few defs for \oe \OE \ae \AE \o \O \ss... (just that!)
%        probably when modifying the code for LEX and 2e).        --bg 03/07/95
%       TeX3+CMR warning message rewritten in english and no more
%        issued twice.                                            --bg 03/08/95
%%
%%      checksum        = "51431 1003 3008 50021"
%%
\def\@tempc{\let\ifLEX\iffalse\let\ifLEXhic\iftrue}%
\ifx\ifLEX\undefined\message{keyboard,}\@tempc\fi%
\let\@tempc=\undefined% so if someone use it...
%%%%%%%%%%%%%%%%%%%%%%%%%%%%%%%%%%%%%%%%%%%%%%%%%%%%%%%%%%%%%%%%%%%%%%%%%
% This code is able to read keyboard.dat and do appropriate customisation
%%%%%%%%%%%%%%%%%%%%%%%%%%%%%%%%%%%%%%%%%%%%%%%%%%%%%%%%%%%%%%%%%%%%%%%%%
%\let\charsubdef\undefined% make MlTeX 3 unkown for tests
%\let\starthyph\undefined%  make MlTeX 2 unkown for tests
%\font\rm=dcr10\input dc-plain.ch% do test with dc fonts
%------------------------------------------------------------------
%\def\tracingall{\tracingonline\@ne\tracingcommands\tw@\tracingstats\tw@
%  \tracingpages\@ne\tracingoutput\@ne\tracinglostchars\@ne
%  \tracingmacros\tw@\tracingparagraphs\@ne\tracingrestores\@ne
%  \showboxbreadth\maxdimen\showboxdepth\maxdimen\errorstopmode}
%
%%%% Dirty hack to be remove later...
\def\@tempc{\let\ifECM\iffalse}%
\ifx\ifECM\undefined\relax\@tempc\fi%
\ifx\documentclass\undefined\else%
  \edef\@tempb{\encodingdefault}\edef\@tempc{OT1}%
    \ifx\dumpORI\undefined\let\dumpORI=\dump%
\wlog{french-kbconfig: Intercepting the dump control sequence for %
                       eliminating few invalid chars}%
          \let\SetCfgFor\undefined%
          \def\dump{\makeatletter\input kbconfig.tex %
\message{french-kbconfig: Writing latex209.cfg ************************}%
                                     \SetCfgFor209%
\message{french-kbconfig: PLEASE COPY latex209.cfg in your LaTeX inputs}%
                                     \ifECM\relax\else\uccode`\^^Y=`\^^Y%
                                                      \uccode`\^^Z=`\^^Z%
\wlog{french-kbconfig: Redefining \string ^\string ^\string Y uccode}%
\wlog{french-kbconfig: Redefining \string ^\string ^\string Z uccode}%
                                     \fi%
                                     \makeatother\dumpORI}%
    \fi%
\fi%
%%%%
\ifx\I\undefined\def\I{I}\fi%
\ifx\SS\undefined\def\SS{\hbox{SS}}\fi% new defs
% \" is redefined, look at the accentuation macros
\ifLEX\else% don't activate for kb2lex
\expandafter\ifx\csname ^^63\endcsname\c % 128-255 going \active locally in V3
            \catcode`^^80=13\catcode`^^81=13\catcode`^^82=13\catcode`^^83=13
            \catcode`^^90=13\catcode`^^91=13\catcode`^^92=13\catcode`^^93=13
            \catcode`^^a0=13\catcode`^^a1=13\catcode`^^a2=13\catcode`^^a3=13
            \catcode`^^b0=13\catcode`^^b1=13\catcode`^^b2=13\catcode`^^b3=13
            \catcode`^^c0=13\catcode`^^c1=13\catcode`^^c2=13\catcode`^^c3=13
            \catcode`^^d0=13\catcode`^^d1=13\catcode`^^d2=13\catcode`^^d3=13
            \catcode`^^e0=13\catcode`^^e1=13\catcode`^^e2=13\catcode`^^e3=13
            \catcode`^^f0=13\catcode`^^f1=13\catcode`^^f2=13\catcode`^^f3=13
            \catcode`^^84=13\catcode`^^85=13\catcode`^^86=13\catcode`^^87=13
            \catcode`^^94=13\catcode`^^95=13\catcode`^^96=13\catcode`^^97=13
            \catcode`^^a4=13\catcode`^^a5=13\catcode`^^a6=13\catcode`^^a7=13
            \catcode`^^b4=13\catcode`^^b5=13\catcode`^^b6=13\catcode`^^b7=13
            \catcode`^^c4=13\catcode`^^c5=13\catcode`^^c6=13\catcode`^^c7=13
            \catcode`^^d4=13\catcode`^^d5=13\catcode`^^d6=13\catcode`^^d7=13
            \catcode`^^e4=13\catcode`^^e5=13\catcode`^^e6=13\catcode`^^e7=13
            \catcode`^^f4=13\catcode`^^f5=13\catcode`^^f6=13\catcode`^^f7=13
            \catcode`^^88=13\catcode`^^89=13\catcode`^^8a=13\catcode`^^8b=13
            \catcode`^^98=13\catcode`^^99=13\catcode`^^9a=13\catcode`^^9b=13
            \catcode`^^a8=13\catcode`^^a9=13\catcode`^^aa=13\catcode`^^ab=13
            \catcode`^^b8=13\catcode`^^b9=13\catcode`^^ba=13\catcode`^^bb=13
            \catcode`^^c8=13\catcode`^^c9=13\catcode`^^ca=13\catcode`^^cb=13
            \catcode`^^d8=13\catcode`^^d9=13\catcode`^^da=13\catcode`^^db=13
            \catcode`^^e8=13\catcode`^^e9=13\catcode`^^ea=13\catcode`^^eb=13
            \catcode`^^f8=13\catcode`^^f9=13\catcode`^^fa=13\catcode`^^fb=13
            \catcode`^^8c=13\catcode`^^8d=13\catcode`^^8e=13\catcode`^^8f=13
            \catcode`^^9c=13\catcode`^^9d=13\catcode`^^9e=13\catcode`^^9f=13
            \catcode`^^ac=13\catcode`^^ad=13\catcode`^^ae=13\catcode`^^af=13
            \catcode`^^bc=13\catcode`^^bd=13\catcode`^^be=13\catcode`^^bf=13
            \catcode`^^cc=13\catcode`^^cd=13\catcode`^^ce=13\catcode`^^cf=13
            \catcode`^^dc=13\catcode`^^dd=13\catcode`^^de=13\catcode`^^df=13
            \catcode`^^ec=13\catcode`^^ed=13\catcode`^^ee=13\catcode`^^ef=13
            \catcode`^^fc=13\catcode`^^fd=13\catcode`^^fe=13\catcode`^^ff=13
\let ^^9c=\undefined%
\fi
\fi%\ifLEX
%%%%%%%%%%%%%%%%%%%%%%%%%%%%%%%%%%%%%%%%%%%%%%%%%%%
%------------------------------------------------------------------
% code coming from compatible.tex (Aug. 1990) of MJF and Justin Bur
\def\csubinverse #1#2#3{\expandafter\xdef\csname #2#3\endcsname{#1}}%
%           #1 -- extended 8bits character to replace accent+letter
%           #2 -- accent macro invocation for...
%           #3 -- this letter (base character)
%------------------------------------------------------------------
\begingroup% to protect all the marmelade done here
\ifLEX\else\let\protect\noexpand\fi%
\catcode`[=13\catcode`@=11%
\catcode`<=13\catcode`>=13% will allow << and >> in keyboard.dat
%%%%%%%%%
% test if ECM can be used (in fact DCM here)
%
\newif\ifECM% local redefinition for a global expansion.
\global\let\ifECM\ifHB% may be it's already set by \ifHB in lhyphen/plain
\ifx\ifHB\undefined\global\ECMfalse\fi%
\ifECM\else% skip this if already decided
\ifx\encodingdefault\undefined%%%%%%%%%%%%
  \ifECM\else% if false test (again) if it might be an ec/dc font
  \ifx\rm\undefined% may be \rm isn't allready defined (specially with LaTeX)
                   \def\graveDC{\mathaccent"7000 }% so test such char
                   \ifx\graveDC\grave\ECMtrue\else\ECMfalse\fi%
  \else% try on fontname that is a better security
     \def\is@DC#1#2#3\relax{\if d#11\fi \if e#11\fi \if c#22\fi}%
     \def\dcdo{% \dcdo code Version 3 as of 01/27/92 on DCFONT discussion list
               \ifnum 12=0\expandafter\is@DC\fontname\font xx\relax\space%
                          \expandafter\dc@first%
               \else      \expandafter\dc@second%
               \fi}%
          \def\dc@first#1#2{#1}%
          \def\dc@second#1#2{#2}%
  {\rm\dcdo{\global\ECMtrue}{\global\ECMfalse}}% final DC test
  \fi%\ifx\rm\undefined
  \fi%\ifECM\else..
%
  \ifECM\global\def\encodingdefault{T1}%   T1 is for DC/EC
   \else\global\def\encodingdefault{OT1}% OT1 is for CMR
  \fi%
%
\else%
\ifx\f@encoding\undefined\edef\f@encoding{\encodingdefault}\fi%
     \edef\@tempa{T1}\ifx\@tempa\f@encoding%
                          {\global\ECMtrue}%
                     \else{\global\ECMfalse}%
                     \fi%
\fi%\encodingdefault
\fi% ifECM
%%%%%%
% Prevent LaTeX2e invalid chars
\ifx\documentclass\undefined\else
\ifECM%
\catcode`\^^_=11%
\catcode`\^^Y=11%
\else%
\catcode`\^^P=11%
\catcode`\^^Q=11%
\catcode`\^^Y=11%
\catcode`\^^Z=11%
\catcode`\^^]=11%
\catcode`\^^[=11%
\catcode`\^^^=11%
\catcode`\^^\=11%
\catcode`\^^_=11%
\fi\fi%
%%%%%%
\newif\ifMACRO\MACROfalse\newif\ifLIGS\LIGSfalse
%------------------------------------------------------------------
% Now where is the accent in the font? local mods of plain code.
% and how code it: diacritic+letter or accented 8bits char?
%
%Caution: never use TeX V2 hexadecimal notation, lowercase would be wrong
\ifECM\def\POSgr{"00}\def\POSac{"01}\def\POSan{"06}
      \def\POSv{"07}\def\POSu{"08}\def\POSeq{"09}
      \def\POShat{"02}\def\POSdot{"0a}\def\POSH{"05}\def\POStil{"03}
      \def\POSum{"04}\def\POSc{"13}
\else\def\POSgr{18}\def\POSac{19}
     \def\POSv{20}\def\POSu{21}\def\POSeq{22}\def\POSan{23}
     \def\POShat{94}\def\POSdot{95}\def\POSH{125}\def\POStil{126}
     \def\POSum{127}\def\POSc{24}
\fi%
\ifLEX\begingroup\catcode`+=0\catcode`[= 1\catcode`]=2
                 \catcode`\{=12\catcode`\}=12\catcode`\\=12
   +gdef+ToEight[% Due to lex capacities we keep only one expansion of
                 % usual accent macros e.g. {\'E} but not \'{E} nor \'E
      +def+`##1[%\\`\{##1\}|\\`##1|
                \{\\`##1\}]
      +def+'##1[%\\'\{##1\}|\\'##1|
                \{\\'##1\}]
      +def+r##1[\\"r "##1|\\r\{##1\}]
      +def+v##1[\\"v "##1|\\v\{##1\}]
      +def+u##1[\\"u "##1|\\u\{##1\}]
      +def+=##1[%\\=\{##1\}|\\=##1|
                \{\\=##1\}]
      +def+^##1[%\\^\{##1\}|\\^##1|
                \{\\^##1\}]
      +def+.##1[%\\.\{##1\}|\\.##1|
                \{\\.##1\}]
      +def+H##1[\\"H "##1|\\H\{##1\}]
      +def+~##1[%\\~\{##1\}|\\~##1|
                \{\\~##1\}]
      +def+"##1[%\\\"\{##1\}|\\\"##1|
                \{\\\"##1\}]
      +def+c##1[\\"c "##1|\\c\{##1\}]
      +def+i[\\i]
      +def+j[\\j]
      +def+copyright[\\"copyright "|\{\\"copyright"\}|\\"copyright"\{\}]
      +def+oe[\\"oe "|\{\\"oe"\}|\\"oe"\{\}]
      +def+OE[\\"OE "|\{\\"OE"\}|\\"OE"\{\}]
      +def+ss[\\"ss "|\{\\"ss"\}|\\"ss"\{\}]
      +def+ae[\\"ae "|\{\\"ae"\}|\\"ae"\{\}]
      +def+AE[\\"AE "|\{\\"AE"\}|\\"AE"\{\}]
      +def+l[\\"l "|\{\\l\}|\\l\{\}]
      +def+L[\\"L "|\{\\L\}|\\L\{\}]
      +def+o[\\"o "|\{\\o\}|\\o\{\}]
      +def+O[\\"O "|\{\\O\}|\\O\{\}]
      +def+P[\\"P "|\{\\P\}|\\P\{\}]
      +def+pounds[\\"pounds "|\{\\"pounds"\}|\\"pounds"\{\}]
      +def+S[\\"S "|\{\\S\}|\\S\{\}]
      +def+fguillemets["+string<+string<"]
      +def+endfguillemets["+string>+string>"]%
            ]%
   +gdef+ToSeven[%
      +def+`##1[{\\`##1}]
      +def+'##1[{\\'##1}]
      +def+r##1[\\r{##1}]
      +def+v##1[\\v{##1}]
      +def+u##1[\\u{##1}]
      +def+=##1[{\\=##1}]
      +def+^##1[{\\^##1}]
      +def+.##1[{\\.##1}]
      +def+H##1[\\H{##1}]
      +def+~##1[{\\~##1}]
      +def+"##1[{\\\"##1}]
      +def+c##1[\\c{##1}]
      +def+i[\\i]
      +def+j[\\j]
      +def+copyright[{\\copyright}]
      +def+oe[{\\oe}]
      +def+OE[{\\OE}]
      +def+ss[{\\ss}]
      +def+ae[{\\ae}]
      +def+AE[{\\AE}]
      +def+l[{\\l}]
      +def+L[{\\L}]
      +def+o[{\\o}]
      +def+O[{\\O}]
      +def+P[{\\P}]
      +def+pounds[{\\pounds}]
      +def+S[{\\S}]
      +def+fguillemets[+string<+string<]
      +def+endfguillemets[+string>+string>]%
            ]%
      +endgroup
\else%
% Folowing \ifx are not required there are here in case of any misplacement
% of kbconfig (eg after hyconfig).
\ifx\documentclass\undefined% do it in 2.09
\edef\`#1{\expandafter\ifx\csname @gr@#1\endcsname\relax%
                           \ifECM{\accent"00 #1}\else{\accent18 #1}\fi%
                       \else\csname @gr@#1\endcsname\fi}%
\edef\'#1{\expandafter\ifx\csname @ac@#1\endcsname\relax%
                           \ifECM{\accent"01 #1}\else{\accent19 #1}\fi%
                       \else\csname @ac@#1\endcsname\fi}%
\edef\v#1{\expandafter\ifx\csname @v@#1\endcsname\relax%
                           \ifECM{\accent"07 #1}\else{\accent20 #1}\fi%
                       \else\csname @v@#1\endcsname\fi}%
\edef\u#1{\expandafter\ifx\csname @u@#1\endcsname\relax%
                           \ifECM{\accent"08 #1}\else{\accent21 #1}\fi%
                       \else\csname @u@#1\endcsname\fi}%
\edef\=#1{\expandafter\ifx\csname @eq@#1\endcsname\relax%
                           \ifECM{\accent"09 #1}\else{\accent22 #1}\fi%
                       \else\csname @eq@#1\endcsname\fi}%
\edef\^#1{\expandafter\ifx\csname @hat@#1\endcsname\relax%
                           \ifECM{\accent"02 #1}\else{\accent94 #1}\fi%
                       \else\csname @hat@#1\endcsname\fi}%
\edef\.#1{\expandafter\ifx\csname @dot@#1\endcsname\relax%
                           \ifECM{\accent"0a #1}\else{\accent95 #1}\fi%
                       \else\csname @dot@#1\endcsname\fi}%
\edef\H#1{\expandafter\ifx\csname @H@#1\endcsname\relax%
                           \ifECM{\accent"05 #1}\else{\accent125 #1}\fi%
                       \else\csname @H@#1\endcsname\fi}%
\edef\~#1{\expandafter\ifx\csname @til@#1\endcsname\relax%
                           \ifECM{\accent"03 #1}\else{\accent126 #1}\fi%
                       \else\csname @til@#1\endcsname\fi}%
% plain defined it ``wrongly'' in hex ("7F) and so could be wrongly lowercased
\edef\"#1{\expandafter\ifx\csname @um@#1\endcsname\relax%plain define it in hex
                           \ifECM{\accent"04 #1}\else{\accent127 #1}\fi%
                       \else\csname @um@#1\endcsname\fi}%
\def\c#1{\expandafter\ifx\csname @c@#1\endcsname\relax\noexpand\c#1%
                       \else\csname @c@#1\endcsname\fi}%
% D. Knuth defined \oalign in a way that we can't use it in
%          an \edef or \xdef so i changed it temporarily
\def\oalign#1{\leavevmode\vtop{\baselineskip\z@skip \lineskip.25ex%
  \ialign{####\crcr#1\crcr}}}% put characters over each other
% our new def for angstrom bec. undefined in plain as an accent. macro
\def\r#1{\if#1a\@a:a\else\if#1A\@A:A\fi\fi}%
\def\@a:a{\expandafter\ifx\csname @an@a\endcsname\relax%
                          \ifECM{\accent"06a}\else{\accent23a}\fi%
                       \else\csname @an@a\endcsname\fi}%
\def\@A:A{\expandafter\ifx\csname @an@A\endcsname\relax%
                          \leavevmode\setbox0\hbox{h}\dimen@\ht0%
                            \advance\dimen@-1ex\rlap{\raise.67\dimen@\hbox{%
                                    \ifECM\char"06\else\char'27\fi}}A%
                       \else\csname @an@A\endcsname\fi}%
\xdef\aa{\r a}\xdef\AA{\r A}%
\fi% done for 2.09
\fi%\ifLEX
% Strings def are global to allow \@accent@ to work properly (not for 2e)
\def\@gobbletwo#1#2{}%
\ifx\documentclass\undefined%
\let\Gdef=\gdef%
\else\let\Gdef\@gobbletwo%
%%%%%%
% Prevent LaTeX2e invalid chars even if gobbled
\catcode`\^^P=11%
\catcode`\^^Q=11%
\catcode`\^^Y=11%
\catcode`\^^Z=11%
\catcode`\^^]=11%
\catcode`\^^[=11%
\catcode`\^^^=11%
\catcode`\^^\=11%
\catcode`\^^_=11%
%%%%%%
\fi%
\ifECM\Gdef\oe{\string^^f7}\else\Gdef\oe{^^[}\fi%
\ifECM\Gdef\OE{\string^^d7}\else\Gdef\OE{^^^}\fi%
\ifECM\Gdef\ss{\string^^ff}\else\Gdef\ss{^^Y}\fi%
\ifECM\Gdef\ae{\string^^e6}% ^^e6 introduce a pb, why?
                          \else%
                          \ifx\undefined\starthyph\Gdef\ae{^^Z}% for V3
                              \else\Gdef\ae{\char`^^Z}\fi%       for V2
                          \fi%
\ifECM\Gdef\AE{\string^^c6}\else\Gdef\AE{^^]}\fi%
\ifECM\else\global\lccode`^^]=`^^Z% a bypass for V2
           \global\uccode`^^Z=`^^]%
           \global\uccode`^^[=`^^^%
           \global\lccode`^^^=`^^[%
           \global\lccode`^^_=`^^\%
           \global\uccode`^^\=`^^_%
           \global\uccode`^^P=`I%
           \global\uccode`^^Q=`J%
\fi%
\ifECM\Gdef\o {\string^^f8}\else\Gdef\o {^^\}\fi%
\ifECM\Gdef\O {\string^^d8}\else\Gdef\O {^^_}\fi%
% but require an uppercase def for \ss different for ec/dc ie:
\ifECM\Gdef\SS{\string^^df}\fi% valid global def for DC
%
\def\D{\char^^bf}%
\ifx\D\pounds\Gdef\pounds{\string^^bf}\fi%
% Plain has no definition of \pounds
\ifx\pounds\undefined{\ifECM\Gdef\pounds{\string ^^bf}%lfonts.tex will redef.it
                     \else\Gdef\pounds{{\csname it\endcsname\char'44}}\fi}\fi%
%
\ifx<\undefined\gdef<{\string<}\fi%  defs until french.sty were activated
\ifx>\undefined\gdef>{\string>}\fi%
%
\ifECM\Gdef\i{^^Y}\else\Gdef\i{^^P}\fi% needed by \csubinverse for patterns
\ifECM\Gdef\j{^^Z}\else\Gdef\j{^^Q}\fi% dotless letters
% Localy needed anyway:
\ifECM\def\i{^^Y}\else\def\i{^^P}\fi% needed by \csubinverse for patterns
% LaTeX 2e doesn't care of i and j without dot on it
\ifx\documentclass\undefined\else%
 \ifECM\edef\i{i}\edef\j{j}\fi%
\fi%
%
% Are we in TeX V2 or V3?
\newif\ifVtwo \newif\ifSkip \Skipfalse \newif\ifISO \ISOtrue % ISO 8859-1/Cork
\newif\ifMlTeXthree \MlTeXthreefalse
\catcode`^=\active% non hexa processing in V2
\ifx\undefined\lefthyphenmin\Vtwotrue% Vtwo means MlTeX V2
                          \def^#1{\ifx#1^\else\string^#1\fi}% eat ^^ in TeX v2
                          \ifx\starthyph\undefined%
                            \message{ATTENTION : TeX Version 2 ne permet pas
                                    d'utiliser des caracteres accentues 8bits}%
                            \ifnum\catcode`\{=1%\} emacs
                            \message{Vous devez :}
                            \message{- soit utiliser un moteur MlTeX,}
                            \message{- soit passer en encodage OT1 par defaut.}
                            \Stop% a vous de choisir... --bg
                            \fi%
                              \Skiptrue\fi%
                     \else\ifx\starthyph\undefined\Vtwofalse\else\Vtwotrue\fi%
                          \ifx\charsubdef\undefined\else\MlTeXthreetrue\fi%
\ifx\SetCfgFor\undefined%
\ifLEX\else%
  \ifMlTeXthree\else%
   \ifECM\else%
   {\catcode`\^=12\newlinechar=`\^
    \message{^french-kbconfig ****Warning****: your TeX V3 engine +  CMR}}
   \message{fonts (your format default) isn't sufficient to hyphenate}%
   \message{words containing diacritics (like in French).}%
   \fi%
  \fi%
\fi% ifLEX
\fi%\SerCfg
\fi% lefthyphenmin
\ifVtwo\else\catcode`\^=7% reset value of plain
% reset first 32 chars too erroneously (?) illegal since 2e
\catcode`\^^00=11%
\catcode`\^^01=11%
\catcode`\^^02=11%
\catcode`\^^03=11%
\catcode`\^^04=11%
\catcode`\^^05=11%
\catcode`\^^06=11%
\catcode`\^^07=11%
\catcode`\^^08=11%
%\catcode`\^^09=11% is tab ^^I
%\catcode`\^^0a=11% is nl  ^^J
\catcode`\^^0b=11%
%\catcode`\^^0c=11% is ff  ^^L
%\catcode`\^^0d=11% is cr  ^^M
\catcode`\^^0e=11%
\catcode`\^^0f=11%
\catcode`\^^10=11%
\catcode`\^^11=11%
\catcode`\^^12=11%
\catcode`\^^13=11%
\catcode`\^^14=11%
\catcode`\^^15=11%
\catcode`\^^16=11%
\catcode`\^^17=11%
\catcode`\^^18=11%
\catcode`\^^19=11%
\catcode`\^^1a=11%
\catcode`\^^1b=11%
\catcode`\^^1c=11%
\catcode`\^^1d=11%
\catcode`\^^1e=11%
\catcode`\^^1f=11%
\fi%
\ifSkip\else% The code is only for MlTeX or TeX V3.xxxx
%%%%%%%%%%%%%%%%%%%%%%%%%%%%%%%%%%%%%%%%%%%%%%%%%
% Codage du tableau : definition de l'accentuation
\def\ac[ #1 | #2 | #3 | #4 | #5 | #6 | #7 | #8 ]%
       {\ifLEX\message{^^J PROCESSED CHARACTERS ARE:}\fi
        \ifx.#1\let\ACa\undefined\else\def\ACa{#1}\fi% si different d'un point
        \ifx.#2\let\ACb\undefined\else\def\ACb{#2}\fi%
        \ifx.#3\let\ACc\undefined\else\def\ACc{#3}\fi%
        \ifx.#4\let\ACd\undefined\else\def\ACd{#4}\fi%
        \ifx.#5\let\ACe\undefined\else\def\ACe{#5}\fi%
        \ifx.#6\let\ACf\undefined\else\def\ACf{#6}\fi%
        \ifx.#7\let\ACg\undefined\else\def\ACg{#7}\fi%
        \ifx.#8\let\ACh\undefined\else\def\ACh{#8}\fi\let[\minu%]
       }%
% La macro precisant les noms internes
\def\noms| #1 | #2 | #3 | #4 | #5 | #6 | #7 | #8 |%
       {\ifx.#1\let\ANa\undefined\else\def\ANa{#1}\fi%
        \ifx.#2\let\ANb\undefined\else\def\ANb{#2}\fi%
        \ifx.#3\let\ANc\undefined\else\def\ANc{#3}\fi%
        \ifx.#4\let\ANd\undefined\else\def\ANd{#4}\fi%
        \ifx.#5\let\ANe\undefined\else\def\ANe{#5}\fi%
        \ifx.#6\let\ANf\undefined\else\def\ANf{#6}\fi%
        \ifx.#7\let\ANg\undefined\else\def\ANg{#7}\fi%
        \ifx.#8\let\ANh\undefined\else\def\ANh{#8}\fi}%
\newif\ifONE\newif\ifOK%
\def\ONECHAR#1#2/{\ifx9#2\ONEtrue\else\ONEfalse\fi}%
%definition des minuscules
%==>\madef{caract accentue}{accentuation}{Nom minuscule}{lettre}
\def\madef#1#2#3#4{\ifx.#1\else\@madef{#1}{#2}{#3}{#4}\fi}% si pas un .
\def\@madef#1#2#3#4{\ONECHAR#19/%
    \ifx#2\undefined%
    \else\ifVtwo\ifONE\xdef#1{#2{#4}}\fi% #1 is defaultly \active with MlTeX V2
         \else\def\tempb{\catcode`}% now in V3
                  \let\A=A\ifMACRO\let\A\active\else\ifLIGS\let\A\active\fi\fi%
                  \ifx#2\empty\let\A=\active\fi%
                  \ifcat\active\A%
                     \ifECM\OKfalse% normally nothing to do
                       \ifMACRO\OKtrue\fi% except for macros
                       \ifLIGS\OKtrue\fi% except for ligs
                       \ifISO\ifx#4\ss\OKtrue\fi% \ss is specific
                        \else\OKtrue\fi% except if not ISO
                       \ifOK%
                          \ifONE% exclude multiletters
                                \ifLEX\else\gdef #1{#4}\fi%
                          \edef\OldActiveChars{\@tempc}%
                          \xdef\@tempc{\OldActiveChars%
                                            \tempb\string#1=\active}%
                          \fi%
                       \fi%
                     \else% Not using ECM
                          \ifONE% exclude multiletters
                                \ifMlTeXthree\else\tempb#1=\active\fi%
                             \ifLEX\else%
                                   \gdef #1{#4}% don't xdef for oe, ae, etc.
                             \fi%ifLEX
                                \edef\OldActiveChars{\@tempc}%
                                \xdef\@tempc{\OldActiveChars%
                                                  \tempb\string#1=\active}%
                          \fi%
                     \fi%
                  \else% Not a Macro, Not a lig, ..
%%%                       \ifECM\else%  no accentuation in 8bits ???
                        \ifLEX\else\xdef#1{\expandafter\noexpand#2{#4}}\fi%
                       \ifMlTeXthree\else\tempb#1=\active\fi%
                       \edef\OldActiveChars{\@tempc}%
                       \xdef\@tempc{\OldActiveChars%
                                         \tempb\string#1=\active}%
%%%                       \fi%
                  \fi%
         \fi% ifVtwo
     \expandafter\def\csname #3\endcsname{#1}% for further use
     \fi% ifx#2
 \ifLEX\ifLEXhic\ifONE\message{\string#1 \space}
  {\ToSeven\def\GUILL{"}
              \edef\FULL{#2\LaLettre}%
  \ifx\empty#2\edef\FULL{\LaLettre}\fi
  \immediate\write\kbseven{\string#1\space\space\space
                                     printf(\GUILL\FULL\GUILL);}}
  {\ToEight\def\GUILL{\ifLIGS"\fi}%
              \edef\FULL{#2{\GUILL\LaLettre\GUILL}}%
  \ifx\empty#2\edef\FULL{\GUILL\LaLettre\GUILL}\fi
  \immediate\write\kbeight{\FULL%
                           |\string#1\space\space\space printf("\string#1");}}
 \fi\fi\fi% ifLEX
}% @madef
\def\minu#1 | #2 | #3 | #4 | #5 | #6 | #7 | #8 | #9 |%
    {\ifLEX\edef\@tempc{}\fi% required for LaTeX kb2lex!
\def\@tempa##1/{\def\LaLettre{##1}}\@tempa#1/%
     \madef#2{\ACa}{LAa}{#1}%
     \madef#3{\ACb}{LAb}{#1}%
     \madef#4{\ACc}{LAc}{#1}%
     \madef#5{\ACd}{LAd}{#1}%
     \madef#6{\ACe}{LAe}{#1}%
     \madef#7{\ACf}{LAf}{#1}%
     \madef#8{\ACg}{LAg}{#1}%
     \madef#9{\ACh}{LAh}{#1}%
     \let\A=A\ifMACRO\let\A\active\else\ifLIGS\let\A\active\fi\fi%
     \ifcat\A\active\xdef\Lchar{}\else\xdef\Lchar{#1}\fi%dummy for \oe \ae etc.
     \let\hex\hexminu}%
\def\hexminu| #1 | #2 | #3 | #4 | #5 | #6 | #7 | #8 %[
                                                    ]%
    {\ifx.#1\let\HLa\undefined\else\edef\HLa{\noexpand#1}\fi%
     \ifx.#2\let\HLb\undefined\else\edef\HLb{\noexpand#2}\fi%
     \ifx.#3\let\HLc\undefined\else\edef\HLc{\noexpand#3}\fi%
     \ifx.#4\let\HLd\undefined\else\edef\HLd{\noexpand#4}\fi%
     \ifx.#5\let\HLe\undefined\else\edef\HLe{\noexpand#5}\fi%
     \ifx.#6\let\HLf\undefined\else\edef\HLf{\noexpand#6}\fi%
     \ifx.#7\let\HLg\undefined\else\edef\HLg{\noexpand#7}\fi%
     \ifx.#8\let\HLh\undefined\else\edef\HLh{\noexpand#8}\fi%
     \let[\MAJU%]
    }%
%definition des majuscules
%==>\Madef{caract accentue}{accentuation}{Nom majuscule}{lettre}
\def\Madef#1#2#3#4{\ifx.#1\else\@madef{#1}{#2}{#3}{#4}\fi}% si pas un .
\def\MAJU#1 | #2 | #3 | #4 | #5 | #6 | #7 | #8 | #9 |%
    {\let\MAa\undefined\let\MAb\undefined\let\MAc\undefined\let\MAd\undefined
     \let\MAe\undefined\let\MAf\undefined\let\MAg\undefined\let\MAh\undefined
     \ifLEX\edef\@tempc{}\fi% required for LaTeX kb2lex!
\def\@tempa##1/{\def\LaLettre{##1}}\@tempa#1/%
     \Madef#2{\ACa}{MAa}{#1}%
     \Madef#3{\ACb}{MAb}{#1}%
     \Madef#4{\ACc}{MAc}{#1}%
     \Madef#5{\ACd}{MAd}{#1}%
     \Madef#6{\ACe}{MAe}{#1}%
     \Madef#7{\ACf}{MAf}{#1}%
     \Madef#8{\ACg}{MAg}{#1}%
     \Madef#9{\ACh}{MAh}{#1}%
     \let\A=A\ifMACRO\let\A\active\else\ifLIGS\let\A\active\fi\fi%
     \ifcat\A\active\xdef\Mchar{}\else\xdef\Mchar{#1}\fi%dummy for \OE etc.
     \let\hex\hexMAJU}%
\def\hexMAJU| #1 | #2 | #3 | #4 | #5 | #6 | #7 | #8 %[
                                                    ]%
    {\ifx.#1\let\HMa\undefined\else\edef\HMa{\noexpand#1}\fi%
     \ifx.#2\let\HMb\undefined\else\edef\HMb{\noexpand#2}\fi%
     \ifx.#3\let\HMc\undefined\else\edef\HMc{\noexpand#3}\fi%
     \ifx.#4\let\HMd\undefined\else\edef\HMd{\noexpand#4}\fi%
     \ifx.#5\let\HMe\undefined\else\edef\HMe{\noexpand#5}\fi%
     \ifx.#6\let\HMf\undefined\else\edef\HMf{\noexpand#6}\fi%
     \ifx.#7\let\HMg\undefined\else\edef\HMg{\noexpand#7}\fi%
     \ifx.#8\let\HMh\undefined\else\edef\HMh{\noexpand#8}\fi%
     % process hexdecimal codes if necessary
     \ifx\HLa\undefined\else\DOhex\ACa{\LAa}{\HLa}{\MAa}{\HMa}{\ANa}\fi%
     \ifx\HLb\undefined\else\DOhex\ACb{\LAb}{\HLb}{\MAb}{\HMb}{\ANb}\fi%
     \ifx\HLc\undefined\else\DOhex\ACc{\LAc}{\HLc}{\MAc}{\HMc}{\ANc}\fi%
     \ifx\HLd\undefined\else\DOhex\ACd{\LAd}{\HLd}{\MAd}{\HMd}{\ANd}\fi%
     \ifx\HLe\undefined\else\DOhex\ACe{\LAe}{\HLe}{\MAe}{\HMe}{\ANe}\fi%
     \ifx\HLf\undefined\else\DOhex\ACf{\LAf}{\HLf}{\MAf}{\HMf}{\ANf}\fi%
     \ifx\HLg\undefined\else\DOhex\ACg{\LAg}{\HLg}{\MAg}{\HMg}{\ANg}\fi%
     \ifx\HLh\undefined\else\DOhex\ACh{\LAh}{\HLh}{\MAh}{\HMh}{\ANh}\fi%
     % now ready to start a 4up
     \let[\minu%]
    }%
\def\@tempc{}%
\def\DOhex#1#2#3#4#5#6{%
%                1              2         3           4        5        6
%           accentuation, lettre_min+ code_hex, lettre_maj+code_hex, nom_accent
%look firstly to see if there are more than one letter character
\ONEfalse\ifx\undefined#4\else\expandafter\ONECHAR#49/\fi%
    \ifVtwo% nothing to do in MlTeX V2
           \else% we are with TeX 3 or MlTeX 3
                \ifECM\ifISO\else% as just a chardef adjustement in hex.
                            \ifLEX\else% not for LEX
                            \def\tempa{\catcode`}\expandafter\tempa#2=\active%
                            \expandafter\xdef#2{\expandafter\string#3}%
                            \fi%
                        \ifONE% exclude multiletters
                         \ifx\undefined#4\else%
                            \def\tempa{\catcode`}\expandafter\tempa#4=\active%
                            \expandafter\xdef#4{\expandafter\string#5}%
                         \fi%
                        \fi%
                      \fi%
                \fi% now define \lccode & \uccode %%%%%%%%%%%%%%%%%%%%%%%%%%%%
                \ifx\empty#1\let\X=X%
                         \ifECM\ifISO% don't activate oe ae etc in 8bits
                                     \edef\D{#2}\edef\E{\noexpand\ss}
                                  \ifx \D\E% but let s-tzet
                                  \else%
                             \def\tempa{\catcode`}%\def\tempb{^}%
                             \expandafter\tempa#2=11%
                             \edef\OldActiveChars{\@tempc}%  disactivate
                             \xdef\@tempc{\OldActiveChars%
                                               \tempa\expandafter\string#2=11}%
                             \ifx\undefined#4\else
                                 \ifONE\expandafter\tempa#4=11%
                                 \edef\OldActiveChars{\@tempc}%disactivate
                                 \xdef\@tempc{\OldActiveChars%
                                               \tempa\expandafter\string#4=11}%
                                 \else%
                                   \ifx\undefined#5\else% case hex code is undef
                                     \let\tempc=#5\expandafter\tempa\tempc=11%
                                   \fi% ifx undefined
                                 \fi%
                             \fi
                             \let\X=\active%
                                  \fi% s-tzet
                         \fi\fi%
                \else\let\X\active\ifISO%
                               \def\tempa{\catcode`}\expandafter\tempa#2=11%
                               \ifx\undefined#4\else\expandafter\tempa#4=11\fi%
                                  \fi%
                \ifECM\ifISO\let\X=\active\fi\fi% No lcc/ucc for normal T1 ISO
                \fi%
                \def\tempa{\lccode`}\def\tempb{=`}%
                \ifLIGS\let\X=\active\fi% don't do that for ligs
                \ifMACRO\let\X=\active\fi% don't do that for macros
                \ifcat\X\active% do usual ucc/lcc encoding
                      \global\expandafter\tempa#2\expandafter\tempb#2%
                    \ifx\undefined#4\else
                      \ifONE%
                           \global\expandafter\tempa#4\expandafter\tempb#2%
                           \def\tempa{\global\uccode`}%
                           \expandafter\tempa#2\expandafter\tempb#4%
                           \expandafter\tempa#4\expandafter\tempb#4%
                      \else% do it in hexa for #4 ==> #5 replace #4
                        \ifx\undefined#5\else% hex code may be undef
                        \global\expandafter\tempa #5\expandafter\tempb #2%
                        \ifLEX\else\def\tempa{\global\uccode`}%
                             \edef\D{\expandafter\string#5}\edef\E{#2}%
                             \expandafter\ifx#5\undefined\else\edef\D{#5}\fi%
                         \ifx\D\E%
                              \ifx ^^9c\undefined% do it once
                                 {\def\tempa{\catcode`}\expandafter\tempa#2=11%
                                  \def\tempa{\global\uccode`}%
                                 \expandafter\tempa#2\tempb^^9c}%
                                 \xdef ^^9c{#4}%
                                 \def\tempa{\catcode`}%
                                 \edef\OldActiveChars{\@tempc}%  disactivate
                                 \xdef\@tempc{\OldActiveChars%
                                              \tempa\string ^^9c=13}%
                              \fi%
                          \else%
                              \expandafter\tempa #2\expandafter\tempb #5%
                              \expandafter\tempa #5\expandafter\tempb #5%
                          \fi%
                        \ifECM\ifISO\else%
                        \ifx\D\E\else% case estzet and SS
                \global\expandafter\edef #3{\expandafter\string #3}% dont point
                \global\expandafter\edef #5{\expandafter\string #5}% elsewhere
                        \fi\fi\fi%
                        \fi% ifLEX
                        \fi% ifx undefined
                      \fi%
                    \fi%
                \else\edef\D{#2}% do ucc/lcc encoding*after* expanding macros
                    \global\expandafter\tempa\D\expandafter\tempb\D%
                    \ifx\undefined#4\else\edef\Q{#4}%
                    \global\expandafter\tempa\Q\expandafter\tempb\D%
                    \def\tempa{\uccode`}%
                    \global\expandafter\tempa\D\expandafter\tempb\Q%
                    \global\expandafter\tempa\Q\expandafter\tempb\Q\fi%
                \fi%
                \ifMlTeXthree\OKfalse% MlTeX V3
                    \ifx\Lchar\empty% ............. this is a control sequence
                    \else% ........................ this is a letter+diacritic
                        \ifx\undefined#6\else%maybe we dont have an accent-name
                        \OKtrue\ifLEX\else%
                        \def\tempa{\charsubdef`}\def\tempb{= }%
                        \edef\tempc{`\Lchar}%
                        \global\expandafter\tempa#2%
                        \expandafter\tempb\csname POS#6\endcsname \tempc%
                               \fi%\ifLEX
                        \fi
                    \fi%
                    \ifx\Mchar\empty% and now uppercase
                    \else%
                        \ifx\undefined#6\else%maybe we dont have an accent-name
                        \OKtrue\ifLEX\else%
                        \def\tempa{\charsubdef`}\def\tempb{= }%
                        \edef\tempc{`\Mchar}%
                        \global\expandafter\tempa#4%
                        \expandafter\tempb\csname POS#6\endcsname \tempc%
                               \fi%\ifLEX
                        \fi
                    \fi%
                    \ifOK%
                      \def\tempb{\catcode`}%
                      \edef\OldActiveChars{\@tempc}% in fact disactivate
                      \xdef\@tempc{\OldActiveChars%
                                        \tempb\expandafter\string#2=11}%
                      \ifx\undefined#4\else%
                          \ifONE%
                         \edef\OldActiveChars{\@tempc}% disactivate
                         \xdef\@tempc{\OldActiveChars%
                                            \tempb\expandafter\string#4=11}%
                          \fi%
                      \fi%
                    \fi% ifOK
                \fi% MlTeX3
           \fi%\ifVtwo
                \ifx\undefined#6\else%
                    \csubinverse {\expandafter\string#3}{@#6@}\Lchar%
                    \csubinverse {\expandafter\string#5}{@#6@}\Mchar%
                \fi
                  }%end of DOhex
\ifx\@@input\undefined\let\@@input\input\fi
\@@input keyboard.dat % load my keyboard specs
\fi% de \ifSkip
\endgroup%
\expandafter\ifx\csname ^^63\endcsname\c % 128-255 going chars in V3
            \catcode`^^80=11\catcode`^^81=11\catcode`^^82=11\catcode`^^83=11
            \catcode`^^90=11\catcode`^^91=11\catcode`^^92=11\catcode`^^93=11
            \catcode`^^a0=11\catcode`^^a1=11\catcode`^^a2=11\catcode`^^a3=11
            \catcode`^^b0=11\catcode`^^b1=11\catcode`^^b2=11\catcode`^^b3=11
            \catcode`^^c0=11\catcode`^^c1=11\catcode`^^c2=11\catcode`^^c3=11
            \catcode`^^d0=11\catcode`^^d1=11\catcode`^^d2=11\catcode`^^d3=11
            \catcode`^^e0=11\catcode`^^e1=11\catcode`^^e2=11\catcode`^^e3=11
            \catcode`^^f0=11\catcode`^^f1=11\catcode`^^f2=11\catcode`^^f3=11
            \catcode`^^84=11\catcode`^^85=11\catcode`^^86=11\catcode`^^87=11
            \catcode`^^94=11\catcode`^^95=11\catcode`^^96=11\catcode`^^97=11
            \catcode`^^a4=11\catcode`^^a5=11\catcode`^^a6=11\catcode`^^a7=11
            \catcode`^^b4=11\catcode`^^b5=11\catcode`^^b6=11\catcode`^^b7=11
            \catcode`^^c4=11\catcode`^^c5=11\catcode`^^c6=11\catcode`^^c7=11
            \catcode`^^d4=11\catcode`^^d5=11\catcode`^^d6=11\catcode`^^d7=11
            \catcode`^^e4=11\catcode`^^e5=11\catcode`^^e6=11\catcode`^^e7=11
            \catcode`^^f4=11\catcode`^^f5=11\catcode`^^f6=11\catcode`^^f7=11
            \catcode`^^88=11\catcode`^^89=11\catcode`^^8a=11\catcode`^^8b=11
            \catcode`^^98=11\catcode`^^99=11\catcode`^^9a=11\catcode`^^9b=11
            \catcode`^^a8=11\catcode`^^a9=11\catcode`^^aa=11\catcode`^^ab=11
            \catcode`^^b8=11\catcode`^^b9=11\catcode`^^ba=11\catcode`^^bb=11
            \catcode`^^c8=11\catcode`^^c9=11\catcode`^^ca=11\catcode`^^cb=11
            \catcode`^^d8=11\catcode`^^d9=11\catcode`^^da=11\catcode`^^db=11
            \catcode`^^e8=11\catcode`^^e9=11\catcode`^^ea=11\catcode`^^eb=11
            \catcode`^^f8=11\catcode`^^f9=11\catcode`^^fa=11\catcode`^^fb=11
            \catcode`^^8c=11\catcode`^^8d=11\catcode`^^8e=11\catcode`^^8f=11
%           \catcode`^^9c=11% this is our bypass
                            \catcode`^^9d=11\catcode`^^9e=11\catcode`^^9f=11
            \catcode`^^ac=11\catcode`^^ad=11\catcode`^^ae=11\catcode`^^af=11
            \catcode`^^bc=11\catcode`^^bd=11\catcode`^^be=11\catcode`^^bf=11
            \catcode`^^cc=11\catcode`^^cd=11\catcode`^^ce=11\catcode`^^cf=11
            \catcode`^^dc=11\catcode`^^dd=11\catcode`^^de=11\catcode`^^df=11
            \catcode`^^ec=11\catcode`^^ed=11\catcode`^^ee=11\catcode`^^ef=11
            \catcode`^^fc=11\catcode`^^fd=11\catcode`^^fe=11\catcode`^^ff=11
%%%%%%
% Prevent LaTeX2e invalid chars
\ifx\documentclass\undefined\else\ifECM
\catcode`\^^_=11%
\catcode`\^^Y=11%
\catcode`\^^Z=11%
\catcode`\^^P=11%
\catcode`\^^Q=11%
\else%
\catcode`\^^P=11%
\catcode`\^^Q=11%
\catcode`\^^Y=11%
\catcode`\^^Z=11%
\catcode`\^^]=11%
\catcode`\^^[=11%
\catcode`\^^^=11%
\catcode`\^^\=11%
\catcode`\^^_=11%
\fi\fi%
%%%%%%
% now intitialize only \ae \ae etc.
\@tempc%
%\show\@tempc%
%%%%%%%% correct bug in latex209.def regarding catcodes %%%%%%%
\edef\SetCfgFor209{{\noexpand\newwrite\noexpand\kbclean%
                    \noexpand\immediate\noexpand\openout\noexpand\kbclean%
                    =latex209.cfg%
                    \noexpand\immediate\noexpand\write\noexpand\kbclean%
                    {\@tempc}%
      \ifECM% Force T1 if we are in T1
                    \noexpand\immediate\noexpand\write\noexpand\kbclean%
                    {\noexpand\noexpand\noexpand\def%
                     \noexpand\noexpand\noexpand\encodingdefault{T1}}%
      \else% Force T1 for non-MlTeX
                    \noexpand\immediate\noexpand\write\noexpand\kbclean%
                    {\noexpand\noexpand\noexpand\ifx%
                     \noexpand\noexpand\noexpand\charsubdef%
                     \noexpand\noexpand\noexpand\undefined%
                     \noexpand\noexpand\noexpand\def%
                     \noexpand\noexpand\noexpand\encodingdefault{T1}%
                     \noexpand\noexpand\noexpand\fi%
                    }%
      \fi%
                   }%
                  }%
%%%%%%%%%%%%%%%%%%%%%%%%%%%%%%%%%%%%%%%%%%%%%%%%%%%%%%%%%%%%%%%
% Reset no more useful macros
\let\@tempc\undefined%
\global\let\Lchar\undefined\global\let\Mchar\undefined%
\ifx\fguillemets\undefined\def\fguillemets{<<}\fi%
\ifx\endfguillemets\undefined\def\endfguillemets{>>}\fi%
%
% Now provide 7-bit new accent macros to produce 8-bit chars
%
%\catcode`@=11% normally done
\let\next\undefined
\ifx\starthyph\undefined\else\let\next\relax\fi
\ifx\charsubdef\undefined\else\let\next\relax\fi
\ifECM\let\next\relax\fi
%%
\ifx\documentclass\undefined%
     \let\Edef\edef\let\Def\def\let\Let\let%
\else%
   \def\Edef#1{\expandafter\relax% define only for MlTeX OT1
               \expandafter\global%
               \expandafter\edef\csname OT\string1\string#1\endcsname}
    \def\Def#1{\expandafter\relax% define only for MlTeX OT1
               \expandafter\global%
               \expandafter\def\csname OT\string1\string#1\endcsname}
    \def\Let#1{\expandafter\relax% define only for MlTeX OT1
               \expandafter\global%
               \expandafter\let\csname OT\string1\string#1\endcsname}

\ifECM\Edef\i{^^Y}\else\Edef\i{^^P}\fi% needed by \csubinverse for patterns
\ifECM\Edef\j{^^Z}\else\Edef\j{^^Q}\fi% dotless letters
%\let\next\undefined% to get rid of the following
\fi%
\ifx\next\relax% then do it (we have 8bits chars)
\Edef\`#1{{\noexpand\expandafter\noexpand\ifx%
          \noexpand\csname @gr@#1\noexpand\endcsname\noexpand\relax%
          \ifECM{\accent"00 #1}\else{\accent18 #1}\fi%
                               \noexpand\else%
          \noexpand\csname @gr@#1\noexpand\endcsname\noexpand\fi}}%
\Edef\'#1{{\noexpand\expandafter\noexpand\ifx%
          \noexpand\csname @ac@#1\noexpand\endcsname\noexpand\relax%
          \ifECM{\accent"01 #1}\else{\accent19 #1}\fi%
                               \noexpand\else%
          \noexpand\csname @ac@#1\noexpand\endcsname\noexpand\fi}}%
\Edef\v#1{{\noexpand\expandafter\noexpand\ifx%
          \noexpand\csname @v@#1\noexpand\endcsname\noexpand\relax%
          \ifECM{\accent"07 #1}\else{\accent20 #1}\fi%
                               \noexpand\else%
          \noexpand\csname @v@#1\noexpand\endcsname\noexpand\fi}}%
\Edef\u#1{{\noexpand\expandafter\noexpand\ifx%
          \noexpand\csname @u@#1\noexpand\endcsname\noexpand\relax%
          \ifECM{\accent"08 #1}\else{\accent21 #1}\fi%
                               \noexpand\else%
          \noexpand\csname @u@#1\noexpand\endcsname\noexpand\fi}}%
\Edef\=#1{{\noexpand\expandafter\noexpand\ifx%
          \noexpand\csname @eq@#1\noexpand\endcsname\noexpand\relax%
          \ifECM{\accent"09 #1}\else{\accent22 #1}\fi%
                               \noexpand\else%
          \noexpand\csname @eq@#1\noexpand\endcsname\noexpand\fi}}%
\Edef\^#1{{\noexpand\expandafter\noexpand\ifx%
          \noexpand\csname @hat@#1\noexpand\endcsname\noexpand\relax%
          \ifECM{\accent"02 #1}\else{\accent94 #1}\fi%
                               \noexpand\else%
          \noexpand\csname @hat@#1\noexpand\endcsname\noexpand\fi}}%
\Edef\.#1{{\noexpand\expandafter\noexpand\ifx%
          \noexpand\csname @dot@#1\noexpand\endcsname\noexpand\relax%
          \ifECM{\accent"0a #1}\else{\accent95 #1}\fi%
                               \noexpand\else%
          \noexpand\csname @dot@#1\noexpand\endcsname\noexpand\fi}}%
\Edef\H#1{{\noexpand\expandafter\noexpand\ifx%
          \noexpand\csname @H@#1\noexpand\endcsname\noexpand\relax%
          \ifECM{\accent"05 #1}\else{\accent125 #1}\fi%
                               \noexpand\else%
          \noexpand\csname @H@#1\noexpand\endcsname\noexpand\fi}}%
\Edef\~#1{{\noexpand\expandafter\noexpand\ifx%
          \noexpand\csname @til@#1\noexpand\endcsname\noexpand\relax%
          \ifECM{\accent"03 #1}\else{\accent126 #1}\fi%
                               \noexpand\else%
          \noexpand\csname @til@#1\noexpand\endcsname\noexpand\fi}}%
\Edef\"#1{{\noexpand\expandafter\noexpand\ifx%
          \noexpand\csname @um@#1\noexpand\endcsname\noexpand\relax%
          \ifECM{\accent"04 #1}\else{\accent127 #1}\fi%
                               \noexpand\else%
          \noexpand\csname @um@#1\noexpand\endcsname\noexpand\fi}}%
%\ifx\@c@@\undefined\let\@c@@=\c\fi% now wrong with 2e
\ifx\@c@@\undefined% needed for OT1
\def\@c@@#1{\setbox\z@\hbox{#1}\ifdim\ht\z@=1ex\accent24 #1%
  \else{\ooalign{\unhbox\z@\crcr\hidewidth\char24\hidewidth}}\fi}%
\fi%
\Def\c#1{{\expandafter\ifx%
         \csname @c@#1\endcsname\relax%
         \@c@@{#1}%
                               \else%
         \csname @c@#1\endcsname\fi}}%
\ifx\@an@a\undefined\else\Let\aa\@an@a\fi%
\ifx\@an@A\undefined\else\Let\AA\@an@A\fi%
\ifx\r\undefined%
\global\Def\r#1{\if#1a\@a:a\else\if#1A\@A:A\fi\fi}%
\global\Def\@a:a{\expandafter\ifx\csname @an@a\endcsname\relax%
                          \ifECM{\accent"06a}\else{\accent23a}\fi%
                       \else\csname @an@a\endcsname\fi}%
\global\Def\@A:A{\expandafter\ifx\csname @an@A\endcsname\relax%
                          \leavevmode\setbox0\hbox{h}\dimen@\ht0%
                            \advance\dimen@-1ex\rlap{\raise.67\dimen@\hbox{%
                                    \ifECM\char"06\else\char'27\fi}}A%
                       \else\csname @an@A\endcsname\fi}%
\fi% \r undefined
\fi% we have 8bits chars

\ifx\next\undefined% we don't have any csubsinverse mechanism
%                  % so reset any previous accent macros to usual
\edef\`#1{\ifECM{\accent"00 #1}\else{\accent18 #1}\fi}%
\edef\'#1{\ifECM{\accent"01 #1}\else{\accent19 #1}\fi}%
\edef\v#1{\ifECM{\accent"07 #1}\else{\accent20 #1}\fi}%
\edef\u#1{\ifECM{\accent"08 #1}\else{\accent21 #1}\fi}%
\edef\=#1{\ifECM{\accent"09 #1}\else{\accent22 #1}\fi}%
\edef\^#1{\ifECM{\accent"02 #1}\else{\accent94 #1}\fi}%
\edef\.#1{\ifECM{\accent"0a #1}\else{\accent95 #1}\fi}%
\edef\H#1{\ifECM{\accent"05 #1}\else{\accent125 #1}\fi}%
\edef\~#1{\ifECM{\accent"03 #1}\else{\accent126 #1}\fi}%
\edef\"#1{\ifECM{\accent"04 #1}\else{\accent127 #1}\fi}%
\ifx\@c@@\undefined\else\let\c=\@c@@\fi%
\def\aa{\ifECM{\accent"06a}\else{\accent23a}\fi}%
\def\AA{\leavevmode\setbox0\hbox{h}\dimen@\ht0%
         \advance\dimen@-1ex\rlap{\raise.67\dimen@\hbox{%
         \ifECM\char"06\else\char'27\fi}}A}%
\fi%we don't have any csubsinverse mechanism

\fi
\endinput%%%%%%%%%%%%%%%%%%%%%%%%%%%%%%%%%%%%%%%%%%%%%%%%%%%


